\chapter{Dashboard Website}
\label{ch:05-website}

\section{Peran dan Kerangka Dashboard}

Dashboard adalah satu-satunya antarmuka manusia pada sistem ini: ia menampilkan
keadaan yang sudah tersimpan di basis data dan mengirim perintah ke perangkat
melalui broker. Dashboard tidak berbicara langsung ke perangkat, tidak menyimpan
keadaan di memorinya sendiri, dan tidak menghitung apa pun yang tidak dapat
dihitung ulang dari basis data. Semua pembacaan menuju \code{\DbFile}; semua
penulisan yang memengaruhi perangkat melewati satu fungsi pengiriman perintah.

Halaman dijalankan oleh Streamlit \StreamlitVersion{} (minimum \StreamlitMin)
dengan \NumPages{} halaman yang dipilih dari sidebar. Sebelum halaman mana pun
dirender, kerangka aplikasi menetapkan tiga hal: runtime (broker, worker, sweeper)
diambil dari singleton proses, daftar perangkat dibaca dari basis data, dan
perangkat terpilih disimpan di state sesi. Bila basis data belum berisi satu
perangkat pun, aplikasi berhenti dengan pesan galat alih-alih menampilkan halaman
kosong. \figref{fig:web-flowchart} menunjukkan alur itu, dan
\figref{fig:web-seq-bootstrap} menempatkannya sebagai urutan pesan. Satu kali rerun
merender tepat satu halaman: navigasi Streamlit tidak menjalankan keenam halaman
sekaligus, dan sidebar beserta topbar dirender berurutan dari atas ke bawah sebelum
halaman yang dipilih.

\figureasset[\linewidth]{web-flowchart}{Alur dashboard satu halaman per rerun: sidebar
dan topbar dirender berurutan, lalu satu halaman dipilih dan dirender, berakhir pada
jalur baca atau jalur tulis.}{fig:web-flowchart}

\figureasset[\linewidth]{web-seq-bootstrap}{Urutan penyalaan sesi dashboard: aplikasi membaca
daftar perangkat, pengguna memilih satu perangkat, lalu halaman aktif
dirender.}{fig:web-seq-bootstrap}

\section{Enam Halaman}

\tabref{tab:05-halaman} merinci isi setiap halaman dan tindakan yang menulis ke
basis data. Hanya tiga halaman yang menulis: Alerts (menandai notifikasi dibaca),
Schedules (mengubah jadwal dan mengirimkannya), dan Control (mengirim perintah).

\begin{table}[H]
\centering
\caption{Enam halaman dashboard, isinya, dan aksi tulisnya.}
\label{tab:05-halaman}
\small
\begin{tabular}{@{}p{1.8cm}p{7.1cm}p{6.1cm}@{}}
\toprule
\textbf{Halaman} & \textbf{Isi} & \textbf{Aksi tulis} \\
\midrule
Status & Sisa pakan, jarak sensor, state perangkat, dan penanda dalam jaringan, dari
view \code{v_device_status}. & Tidak ada. \\
Live & Grafik sisa pakan untuk rentang 1, 6, atau 24 jam dan tabel \LiveTail{} pesan
MQTT terakhir dari buffer memori worker. & Tidak ada. \\
Alerts & Notifikasi pakan menipis, pakan habis, pemberian pakan gagal, dan perangkat
offline. & Menandai satu atau semua notifikasi dibaca. \\
History & Riwayat pemberian pakan 1, 7, atau 30 hari beserta level sebelum dan
sesudah, jenis pemicu, dan status. & Tidak ada. \\
Schedules & Daftar jadwal; tambah, ubah, hapus, aktifkan atau nonaktifkan; kirim
jadwal aktif ke perangkat. & Menulis \code{schedules} lalu mengirim
\code{set_schedule}. \\
Control & Pakan manual, permintaan state, ambang notifikasi, penyetelan jam, dan uji
servo. & Menulis \code{commands} dan \code{feeding_logs}; mengubah ambang dan zona
waktu di \code{devices}. \\
\bottomrule
\end{tabular}
\end{table}

\section{Jalur Baca dan Jalur Tulis}

Pembacaan melewati \file{dashboard/data.py} yang membungkus kueri dengan cache
berdurasi pendek, sedangkan penulisan tidak pernah melewati cache: setiap aksi tulis
mengosongkan cache dan memicu render ulang (\tabref{tab:05-cache}).

\begin{table}[H]
\centering
\caption{Lapisan baca dan tulis pada dashboard.}
\label{tab:05-cache}
\small
\begin{tabular}{@{}p{3.2cm}p{5.0cm}p{6.8cm}@{}}
\toprule
\textbf{Lapisan} & \textbf{Mekanisme} & \textbf{Tujuan} \\
\midrule
Baca berulang & Cache dua detik pada seluruh fungsi baca di \file{data.py} &
Menahan agar penyegaran otomatis tidak menjalankan kueri yang sama berkali-kali. \\
Penyegaran otomatis & Fragmen halaman berulang tiap dua detik pada Status, Live, dan
Schedules & Menampilkan data baru tanpa memuat ulang seluruh halaman. \\
Penyegaran manual & Tombol muat ulang mengosongkan cache lalu meminta render ulang &
Memberi pengguna kendali ketika data harus segar seketika. \\
Tulis & Aksi memanggil pengiriman perintah atau kueri tulis, lalu mengosongkan cache
& Menjamin halaman berikutnya membaca hasil tulisan sendiri. \\
\bottomrule
\end{tabular}
\end{table}

State sesi dipakai untuk dua hal saja: kode perangkat terpilih dan pilihan bahasa.
Keduanya per sesi peramban, sehingga dua pengguna tidak saling mengubah pilihan
perangkat. Karena basis data dan broker dipakai bersama, perubahan satu pengguna
terlihat oleh pengguna lain pada penyegaran berikutnya.

\begin{notebox}
Dashboard tidak memiliki autentikasi: siapa pun yang dapat membuka alamatnya dapat
mengirim perintah pakan. Batasan ini konsisten dengan asumsi jaringan lokal pada
\chapref{ch:01-pendahuluan} dan akibatnya dicatat pada
\chapref{ch:09-catatan-asumsi}.
\end{notebox}

\section{Keadaan Perintah di Antarmuka}

\figref{fig:web-state-command} memodelkan apa yang dialami pengguna setelah menekan
tombol pengirim perintah. Ada empat keadaan akhir: sukses, gagal dengan pesan dari
perangkat, tidak pernah dijawab sehingga ditandai \code{timeout}, atau gagal
dikirim sebelum pesan keluar dari platform. Keempatnya terlihat berbeda, dan keadaan
``dikirim'' tidak pernah berarti ``berhasil''. Karena itu halaman Control
menampilkan pengenal perintah setelah tombol ditekan, bukan pesan sukses.

\figureasset[\linewidth]{web-state-command}{Keadaan perintah di antarmuka: perintah disusun,
baris basis data ditulis, pesan dipublikasikan, lalu menunggu ack sampai batas
\CmdTimeoutS{} detik.}{fig:web-state-command}

\section{Alur Fitur}

Lima alur berikut adalah alur yang paling sering dipakai, digambarkan sebagai
urutan pesan antara peramban, aplikasi, basis data, broker, dan perangkat.

\subsection{Pemberian Pakan Manual yang Berhasil}

Pengguna memilih jumlah porsi pada halaman Control dan mengonfirmasi lewat dialog.
Aplikasi menulis baris perintah dan baris riwayat pakan dalam satu transaksi,
mempublikasikan perintah dengan QoS~1, lalu menunggu. Perangkat mengeksekusi porsi
satu per satu, dan setelah porsi terakhir mengukur ulang jarak, mengirim telemetri,
lalu mengirim ack. Aplikasi menandai keduanya sukses dan mengisi level sesudah dari
telemetri terakhir (\figref{fig:web-seq-manual-feed-ok}).

\figureasset[\linewidth]{web-seq-manual-feed-ok}{Urutan pemberian pakan manual yang berhasil,
dari klik tombol sampai baris riwayat pakan berstatus sukses.}{fig:web-seq-manual-feed-ok}

\subsection{Pemberian Pakan Manual yang Gagal}

Kegagalan datang dari perangkat sebagai ack dengan hasil selain \code{ok}, misalnya
\code{error_servo} ketika servo tidak terpasang atau \code{error_low_level} ketika
hopper terbaca kosong. Aplikasi menandai perintah dan riwayat gagal, lalu membuat
notifikasi \code{feed_failed} berkategori kritis
(\figref{fig:web-seq-manual-feed-fail}). Kegagalan berbentuk tidak adanya jawaban
tidak diselesaikan jalur ack, melainkan oleh sweeper setelah \CmdTimeoutS{} detik.

\figureasset[\linewidth]{web-seq-manual-feed-fail}{Urutan pemberian pakan manual yang gagal:
ack berisi kode kesalahan membuat perintah dan riwayat ditandai gagal serta
memunculkan notifikasi.}{fig:web-seq-manual-feed-fail}

\subsection{Menyinkronkan Jadwal ke Perangkat}

Halaman Schedules menyimpan jadwal di basis data lebih dahulu, sehingga perubahan
tidak hilang hanya karena perangkat sedang mati. Pengguna lalu dapat mengirim
seluruh jadwal aktif sebagai satu perintah \code{set_schedule}; perangkat menyimpan
seluruh daftar, menggantikan daftar sebelumnya, dan menjawab dengan jumlah jadwal
yang tersimpan (\figref{fig:web-seq-schedule-sync}).

\figureasset[\linewidth]{web-seq-schedule-sync}{Urutan sinkronisasi jadwal: jadwal aktif dibaca
dari basis data, dikirim sebagai satu perintah, lalu perangkat menjawab dengan
jumlah jadwal yang tersimpan.}{fig:web-seq-schedule-sync}

\subsection{Telemetri dan Peringatan}

Telemetri adalah satu-satunya pesan yang datang tanpa diminta. Setiap pesan
menghasilkan satu baris telemetri, pembaruan waktu terakhir terlihat, penutupan
notifikasi offline, dan evaluasi ambang dengan aturan tunggal: persentase nol
memicu notifikasi pakan habis, persentase di atas nol tetapi tidak lebih dari ambang
memicu notifikasi pakan menipis, dan persentase di atas ambang menutup keduanya.
Sebelum menulis notifikasi baru, aplikasi memeriksa apakah notifikasi sejenis masih
terbuka, sehingga satu kondisi tidak menghasilkan ratusan baris
(\figref{fig:web-seq-telemetry-alert}).

\figureasset[\linewidth]{web-seq-telemetry-alert}{Urutan penanganan telemetri: perhitungan
persentase, penyimpanan baris telemetri, lalu pembukaan atau penutupan notifikasi
sesuai ambang.}{fig:web-seq-telemetry-alert}

\subsection{Backup dan Restore Basis Data}

Panel backup menyalin basis data yang sedang hidup ke berkas sementara lalu
menyajikannya sebagai unduhan. Restore melakukan arah sebaliknya: berkas unggahan
diperiksa apakah benar basis data sistem ini, lalu isinya disalin ke basis data
aktif memakai antarmuka backup SQLite, bukan dengan menukar berkas. Cara itu dipilih
karena koneksi UI dan worker sedang terbuka; menukar berkas di bawah koneksi yang
hidup akan membuat keduanya membaca berkas yang salah. Setelah penyalinan, jurnal
WAL dipadatkan agar berkas utama konsisten (\figref{fig:web-seq-backup-restore}).

\figureasset[\linewidth]{web-seq-backup-restore}{Urutan backup dan restore: salinan konsisten
diunduh sebagai berkas, dan berkas unggahan disalin kembali ke basis data aktif
lewat antarmuka backup SQLite.}{fig:web-seq-backup-restore}

\section{Pemilihan Mode Broker}

\figref{fig:web-flow-broker-mode} memuat alur yang dijalankan saat proses
dinyalakan, yang sudah dibahas dari sisi arsitektur pada \chapref{ch:02-arsitektur}.
Dari sisi dashboard, hasil alur itu menentukan label pada sidebar: broker tertanam,
broker yang sudah berjalan, atau broker eksternal. Label itu satu-satunya cara
pengguna mengetahui mengapa perangkat tidak terlihat, misalnya ketika dashboard
memakai broker eksternal sementara perangkat menyambung ke broker tertanam di mesin
lain.

\figureasset[\linewidth]{web-flow-broker-mode}{Alur pemilihan mode broker saat proses
dinyalakan: pemeriksaan konfigurasi, pemeriksaan port, lalu penyalaan broker
tertanam atau penggunaan broker yang sudah ada.}{fig:web-flow-broker-mode}

\section{Sweeper: Dua Tugas Berkala}

Sweeper bangun tiap \SweepS{} detik dan mengerjakan dua tugas yang tidak saling
bergantung (\tabref{tab:05-sweeper}). Keduanya menyelesaikan keadaan yang tidak akan
pernah diselesaikan oleh pesan masuk, yaitu pesan yang seharusnya datang tetapi
tidak datang.

\begin{table}[H]
\centering
\caption{Dua tugas berkala sweeper.}
\label{tab:05-sweeper}
\small
\begin{tabular}{@{}p{3.2cm}p{4.4cm}p{7.4cm}@{}}
\toprule
\textbf{Tugas} & \textbf{Kriteria} & \textbf{Akibat} \\
\midrule
Penanda perintah kedaluwarsa & Perintah berstatus terkirim yang dibuat lebih dari
\CmdTimeoutS{} detik lalu & Status perintah menjadi \code{timeout}; riwayat pakan
terkait menjadi gagal dengan pesan \code{timeout}. \\
Penanda perangkat offline & Perangkat yang masih ditandai online tetapi tidak
mengirim data lebih dari \OfflineS{} detik & Status perangkat menjadi offline dan
notifikasi offline dibuat bila belum ada. \\
\bottomrule
\end{tabular}
\end{table}

Tugas kedua ada karena broker tertanam tidak menutup sesi client yang diam. Tanpa
penanda itu, perangkat yang mati tanpa memutus soket akan tetap terlihat online
selamanya. Ambang \OfflineS{} detik lebih longgar daripada interval telemetri
\TelemetryMs{} ms, sehingga satu pesan yang terlewat tidak langsung membuat
perangkat dinyatakan mati.

\section{Internasionalisasi}

Seluruh teks yang dilihat pengguna berada di satu modul terpisah dan tersedia dalam
dua bahasa, Indonesia dan Inggris, dengan bahasa bawaan Indonesia. Pilihan bahasa
disimpan di state sesi dan tidak memengaruhi nilai protokol: kode perangkat, topik,
pengenal perintah, dan nilai enum sengaja tidak diterjemahkan karena merupakan
bagian dari kontrak dengan perangkat, bukan teks tampilan.
